% !PRINTED 598
\weberpage{598}
%<<CONT
in turn contains $R$ as a divisor. The numbers of $\Omega'$ remain
unchanged under the substitutions $\Phi'$, and $\Phi'$ is the group of the
field $\Omega$ with respect to $\Omega'$ (cf. Bd. I, \S.\,155). We denote
the degree of $\Phi'$, which is a divisor of $n$, by $n'$, and set

\weq{1}{n = m'\,n'.}

Then $n'$ is the degree of $\Omega$ with respect to $\Omega'$, and $m'$
the degree of $\Omega'$ with respect to $R$.

It is to be observed here that while $\Omega$ is a normal field
also with respect to $\Omega'$, yet $\Omega'$ will in general not be a
normal field with respect to $R$.

The prime ideals in the field $\Omega'$ can now be derived completely
from those of $\Omega$.

If $\mathfrak{p}$ is any prime ideal in $\Omega$, then there is \gesperrt{one and
only one ideal $\mathfrak{p}'$ in $\Omega'$ divisible by $\mathfrak{p}$}, and $\mathfrak{p}'$, as a
divisor of $\mathfrak{P}$, can be divisible by no ideals other than those
conjugate to $\mathfrak{p}$. If $\varphi'$ runs through the substitutions of $\Phi'$, then $\mathfrak{p}'$,
since it remains unchanged under $\varphi'$, is also divisible by $\mathfrak{p}\,|\,\varphi'$.

If therefore $\mathfrak{p}_1'$ is divisible by $\mathfrak{p}\,|\,\varphi_1$, it is also divisible by
$\mathfrak{p}\,|\,\psi\,\varphi_1\,\varphi'$, if $\psi$ denotes an arbitrary element of the group
$\boldsymbol{\Psi}$ (\S.\,162).

If, therefore, $\varphi$ is any element of the system

\[
\Phi_1 = \boldsymbol{\Psi}\,\varphi_1\,\Phi'
\]

then $\mathfrak{p}_1'$ is divisible by $\mathfrak{p}\,|\,\varphi$. Accordingly the
group decomposition that we set forth in \S.\,28 of this Band now
comes into consideration.

If $\varphi_2$ is an element of $\Phi$ not contained in $\Phi_1$, the
system

\[
\Phi_2 = \boldsymbol{\Psi}\,\varphi_2\,\Phi'
\]

has no element at all in common with $\Phi_1$. If there is yet another
element $\varphi_3$, which occurs neither in $\Phi_1$ nor in $\Phi_2$, one forms
in the same way

\[
\Phi_3 = \boldsymbol{\Psi}\,\varphi_3\,\Phi',
\]

and continues in this manner until the whole group $\Phi$ is exhausted. One
thus obtains the decomposition

\weq{2}{\Phi = \Phi_1 + \Phi_2 + \Phi_3 + \cdots + \Phi_{e'}.}

We can determine the degree of one of these systems $\Phi_r$ according
to \S.\,28. It is equal to the product of the degree of $\Phi$ multiplied by % [sic] source prints $\Phi$; the equivalent expression below requires $\boldsymbol\Psi$
the index of the intersection $\Psi_r'$ of $\Phi'$ with $\Psi_r = \varphi_r^{-1}\boldsymbol{\Psi}\varphi_r$ with
respect to $\Phi'$, or, what is the same thing, equal to the product of the
degrees of $\boldsymbol{\Psi}$ and $\Phi'$, divided by the
%CONT>>
